\chapter{Rates Risk}\label{ch:rc:rates-risk}

A dollar swap desk holds two hundred swaps, USD~12 billion of notional, and its
risk report shows a parallel DV01 of zero: a one-basis-point rise of every
rate would leave it exactly where it is. Replay on it the moves of the US
Treasury curve on 21 October 2022, when the two-year yield fell 13 basis
points and the thirty-year rose 9, and it loses three million dollars. The
report was right about the level and silent about the shape. Rates risk is a
vector, one number per point of the curve, and managing it means choosing the
coordinates in which to read the vector, knowing how the curve actually moves,
and hedging with a handful of instruments a risk spread over many. This
chapter does all three on chapter~1's curve builder, with ten years of daily
Treasury data as the model of how curves move.

\section{Bucketed sensitivities and risk ladders}

\begin{definition}[Bucketed sensitivity, risk ladder]\label{def:rc:rates-risk:ladder}
A \emph{bucketed sensitivity}\index{bucketed sensitivity} $\Delta_k$ of a
position is the change in its value when the $k$-th input quote of the curve
rises by one basis point and the curve is recalibrated, the others held. The
\emph{risk ladder}\index{risk ladder} is the vector $(\Delta_1,\dots,\Delta_n)$
over all inputs; its sum is the sensitivity to a parallel shift of the quotes,
and its dot product with a vector of quote moves $\delta q$ (in basis points)
is the first-order P\&L, $\pnl\approx\sum_k\Delta_k\,\delta q_k$.
\end{definition}

A par swap maturing on a \omterm{def:rc:curve-construction:pillar}{pillar} loads only that \omterm{def:rc:curve-construction:pillar}{pillar} under a local
interpolation (\cref{prop:rc:curve-construction:triangular}): on the chapter's
\omterm{def:rc:curve-construction:pillar}{eight-pillar} curve (flat forwards, par SOFR swaps at 1, 2, 3, 5, 7, 10, 20 and
30 years) a ten-year payer of USD~100 million has a ladder of USD~83\,778 on the
ten-year quote and zero elsewhere. A book's ladder is the sum of its trades'.

\begin{example}[The desk's ladder]\label{ex:rc:rates-risk:ladder}
The chapter's book holds 200 swaps of random maturity, direction and size
(seeded), plus a two-year payer of USD~1.471 billion and a thirty-year payer of
USD~281 million added so that the parallel DV01 is zero. Its ladder
(\cref{fig:rc:rates-risk:ladder}) is anything but zero: $+212\,010$ at two years,
$+249\,817$ at seven, $-279\,791$ at ten, $-1\,135\,816$ at twenty and
$+958\,308$ at thirty, per basis point.
\end{example}

\begin{omfigure}
\begin{tikzpicture}
\begin{axis}[width=12cm,height=5.4cm,ybar,bar width=9pt,
  ylabel={USD thousand per bp},
  tick label style={font=\small},label style={font=\small},
  xmin=-0.6,xmax=7.6,ymin=-1250,ymax=1100,grid=major,grid style={gray!25},
  xtick=data,xticklabels from table={figdata/rates-credit-risk/03-rates-risk/ladder.csv}{tenor},
  legend columns=2,legend style={font=\footnotesize,at={(0.5,-0.2)},anchor=north,draw=none,fill=none,/tikz/every even column/.append style={column sep=8pt}},
  table/col sep=comma]
\addplot[fill=omDef!55,draw=omDef] table[x=i,y=book]{figdata/rates-credit-risk/03-rates-risk/ladder.csv};
\addplot[fill=omThm!55,draw=omThm] table[x=i,y=hedged]{figdata/rates-credit-risk/03-rates-risk/ladder.csv};
\legend{book,book with the minimum-variance hedge}
\end{axis}
\end{tikzpicture}

\omcaption{The \omterm{def:rc:rates-risk:ladder}{risk ladder} of the chapter's 200-swap book, whose buckets sum to
zero, and the ladder after the three-swap hedge of
\cref{ex:rc:rates-risk:hedge} (two-, twenty- and thirty-year swaps). The hedge
does not zero those three buckets; it leaves the combination of risks that
historical moves make cheapest to carry. Data: the chapter's illustrative book
and tutorial.}\label{fig:rc:rates-risk:ladder}
\end{omfigure}

\section{From par risk to zero risk: the Jacobian}

The ladder above is in the coordinates of the input quotes. Models, scenario
engines and regulators often want it in other coordinates: zero rates at the
\omterm{def:rc:curve-construction:pillar}{pillars}, forward rates, or a standard set of key maturities.

\begin{definition}[Curve Jacobian]\label{def:rc:rates-risk:jacobian}
The \emph{curve Jacobian}\index{curve Jacobian} is the matrix $J_{kj} =
\partial m_k/\partial z_j$ of the derivatives of the model quotes of the input
instruments with respect to the curve's \omterm{def:rc:curve-construction:pillar}{pillar} parameters, at the calibrated
curve; its inverse gives how the \omterm{def:rc:curve-construction:pillar}{pillars} move when the quotes move, $\delta z
= J^{-1}\delta q$.
\end{definition}

\begin{proposition}[Changing coordinates]\label{prop:rc:rates-risk:jacobian}
Let $g_q = (\partial V/\partial q_k)_k$ and $g_z = (\partial V/\partial
z_j)_j$ be the ladders of a position in quotes and in \omterm{def:rc:curve-construction:pillar}{pillar} zero rates. Then
\[
g_z = J^\top g_q,\qquad g_q = (J^\top)^{-1}g_z .
\]
\end{proposition}
\begin{proof}
The position depends on the quotes only through the \omterm{def:rc:curve-construction:pillar}{pillars}, $V = V(z(q))$, so
$g_q = (\partial z/\partial q)^\top g_z = (J^{-1})^\top g_z$.
\end{proof}

\begin{omfigure}
\begin{tikzpicture}[font=\small,
  box/.style={draw,thick,rounded corners=2pt,align=center,minimum height=1cm,minimum width=2.6cm,inner sep=4pt},
  arr/.style={-{Stealth[length=2.5mm]},thick}]
\node[box,draw=omDef,fill=omDef!8] (q) at (0,0) {input quotes\\$q_1,\dots,q_n$};
\node[box,draw=omProp,fill=omProp!8] (z) at (5,0) {pillar zero rates\\$z_1,\dots,z_n$};
\node[box,draw=omThm,fill=omThm!8] (v) at (10,0) {value of the\\position $V$};
\draw[arr] (q) -- node[above=4pt] {calibration} node[below=4pt] {$\delta z = J^{-1}\delta q$} (z);
\draw[arr] (z) -- node[above=4pt] {pricing} node[below=4pt] {$g_z=\partial V/\partial z$} (v);
\draw[arr,omExo,dashed] (v.south) to[bend left=22] node[below=3pt] {risk in quotes: $g_q = (J^\top)^{-1}g_z$} (q.south);
\end{tikzpicture}

\omcaption{Two coordinate systems for one risk. Calibration maps quotes to
\omterm{def:rc:curve-construction:pillar}{pillars}, pricing maps \omterm{def:rc:curve-construction:pillar}{pillars} to value; the Jacobian of the calibration carries
a ladder from one system to the other
(\cref{prop:rc:rates-risk:jacobian}).}\label{fig:rc:rates-risk:jacobian}
\end{omfigure}

\begin{example}[Two ladders of one swap]\label{ex:rc:rates-risk:twoladders}
The ten-year payer's par ladder is $83\,778$ on the ten-year quote alone; its
zero ladder is $367$, $713$, $1\,531$, $3\,205$, $5\,120$ and $74\,016$ on the
one- to ten-year \omterm{def:rc:curve-construction:pillar}{pillars}. Raising an early zero rate lowers the value of the
fixed coupons the payer owes, and in quote space that effect is absorbed by
recalibration. For the whole book, the par ladder sums to zero and the zero
ladder to USD~11\,408: a parallel shift of quotes and a parallel shift of zero
rates are different scenarios.
\end{example}

\begin{definition}[Key-rate duration]\label{def:rc:rates-risk:krd}
A \emph{key-rate duration}\index{key-rate duration} (Ho, 1992) is the
sensitivity of a position to a ``tent'' shift of the zero curve centred on one
key maturity: one basis point at the key, falling linearly to zero at the two
neighbouring keys. The tents sum to one at every maturity, so key-rate
sensitivities add up to the sensitivity to a parallel shift of zero rates;
they are independent of how the curve was built, which makes them the usual
currency of risk reports across desks.
\end{definition}

\section{Principal components of the curve}

How do curves move? The covariance matrix of daily changes answers, and its
eigenvectors, the principal components of One Quant Book 4, chapter 22, give
the answer a shape.

\begin{definition}[Level, slope and curvature factors]\label{def:rc:rates-risk:factors}
For a curve sampled at $n$ maturities, let $v_1, v_2, v_3$ be the first three
eigenvectors of the covariance of its daily changes (with $\lambda_1\ge
\lambda_2\ge\lambda_3$). The \emph{level factor}\index{level factor} $v_1$ has
loadings of one sign, the \emph{slope factor}\index{slope factor} $v_2$
changes sign once, from short to long maturities, and the \emph{curvature
factor}\index{curvature factor} $v_3$ changes sign twice. The move of day $t$ has
score $v_i^\top\delta q_t$ on factor $i$, and the factors are uncorrelated by
construction.
\end{definition}

\begin{example}[Ten years of Treasury moves]\label{ex:rc:rates-risk:pca}
Over the 2\,681 daily changes of the US Treasury par curve from January 2016
to September 2026 (eight maturities from one to thirty years), the first three
components explain 85.1\%, 11.0\% and 2.1\% of the variance, 98.2\% together,
with daily standard deviations of 13.6, 4.9 and 2.1 basis points of score
(\cref{fig:rc:rates-risk:pca}). The level loads less on the one-year than on
the belly; the \omterm{def:rc:rates-risk:factors}{curvature factor} is dominated by the one-year yield, the
maturity most tied to the next policy meetings. Litterman and Scheinkman found
the same three factors in 1991.
\end{example}

\begin{omfigure}
\begin{tikzpicture}
\begin{axis}[width=11.5cm,height=5.6cm,
  xlabel={maturity (years, log scale)},ylabel={loading},xmode=log,
  xtick={1,2,3,5,7,10,20,30},xticklabels={1,2,3,5,7,10,20,30},
  tick label style={font=\small},label style={font=\small},
  xmin=0.9,xmax=33,ymin=-0.6,ymax=0.9,grid=major,grid style={gray!25},
  legend columns=3,legend style={font=\footnotesize,at={(0.5,-0.25)},anchor=north,draw=none,fill=none,/tikz/every even column/.append style={column sep=8pt}},
  table/col sep=comma]
\addplot[omDef,very thick,mark=*] table[x=m,y=pc1]{figdata/rates-credit-risk/03-rates-risk/pca.csv};
\addplot[omThm,thick,mark=square*] table[x=m,y=pc2]{figdata/rates-credit-risk/03-rates-risk/pca.csv};
\addplot[omProp,thick,dashed,mark=triangle*] table[x=m,y=pc3]{figdata/rates-credit-risk/03-rates-risk/pca.csv};
\legend{level (85.1\%),slope (11.0\%),curvature (2.1\%)}
\end{axis}
\end{tikzpicture}

\omcaption{Loadings of the first three principal components of daily changes of
US Treasury par yields, January 2016 to September 2026, with their shares of
variance. Source: US Treasury, daily par yield curve rates (public domain);
the chapter's tutorial.}\label{fig:rc:rates-risk:pca}
\end{omfigure}

\begin{remark}[The loss of the hook, factor by factor]\label{rem:rc:rates-risk:decomp}
Decomposing the move of 21 October 2022 and the book's ladder on all eight
components, the \omterm{def:rc:rates-risk:factors}{slope factor} (score $+21.2$) accounts for $-\text{USD}~3.09$
million of the $-3.00$ million first-order loss, and the higher components
add $+0.72$, $+0.69$ and $-1.48$ million that nearly cancel. The book was
level-neutral and short a steepening.
\end{remark}

\section{Hedging a swap book}

A book's ladder has as many buckets as the curve has \omterm{def:rc:curve-construction:pillar}{pillars}; a desk hedges it
with a few liquid swaps. Which ones, and how much of each, depends on how the
buckets move together.

\begin{method}[Minimum-variance hedge]\label{met:rc:rates-risk:mvhedge}
Let $g$ be the book's ladder, $H$ the matrix whose columns are the ladders of
one unit of each hedge instrument, and $\Sigma$ the covariance of daily quote
moves. The daily P\&L of book plus hedge is $(g+Hh)^\top\delta q$, of variance
$(g+Hh)^\top\Sigma(g+Hh)$, minimised by
\[
h^\star = -\bigl(H^\top\Sigma H\bigr)^{-1}H^\top\Sigma g .
\]
Report the residual share of variance $(g+Hh^\star)^\top\Sigma(g+Hh^\star) /
g^\top\Sigma g$, and choose the instruments that make it smallest.
\end{method}

\omcode{code/firm/ratesrisk/firm_ratesrisk.py}{94}{101}{The minimum-variance hedge and the variance it leaves.}\label{lst:rc:rates-risk:hedge}

\begin{example}[Three swaps against two hundred]\label{ex:rc:rates-risk:hedge}
With Treasury par-yield changes standing in for swap-rate changes, the book's
daily P\&L has a standard deviation of USD~1\,288\,843. The textbook trio of two-,
ten- and thirty-year swaps leaves 58.3\% of the variance; hedging the three
principal components exactly leaves 68.9\%, because most of the book's risk
lies in the twenty-against-thirty-year spread, which the first three factors do
not describe. The best three swaps are the two-, twenty- and thirty-year:
receive fixed on USD~1.196 billion of two-year, pay on USD~801 million of
twenty-year and receive on USD~515 million of thirty-year. They leave 4.2\%
of the variance, a daily standard deviation of USD~264\,554
(\cref{fig:rc:rates-risk:pnl}).
\end{example}

\begin{omfigure}
\begin{tikzpicture}
\begin{axis}[width=12cm,height=5.4cm,
  xlabel={date},ylabel={daily P\&L (USD million)},
  xtick={2025,2025.5,2026,2026.5},xticklabels={Jan 2025,Jul 2025,Jan 2026,Jul 2026},
  tick label style={font=\small},label style={font=\small},
  xmin=2024.66,xmax=2026.75,ymin=-6,ymax=4.5,grid=major,grid style={gray!25},
  legend columns=2,legend style={font=\footnotesize,at={(0.5,-0.25)},anchor=north,draw=none,fill=none,/tikz/every even column/.append style={column sep=8pt}},
  table/col sep=comma]
\addplot[omDef!70,thin] table[x=t,y=book]{figdata/rates-credit-risk/03-rates-risk/pnl.csv};
\addplot[omThm,thick] table[x=t,y=hedged]{figdata/rates-credit-risk/03-rates-risk/pnl.csv};
\legend{book (first order),book with the three-swap hedge}
\end{axis}
\end{tikzpicture}

\omcaption{First-order daily P\&L of the book, $g^\top\delta q_t$, on the
Treasury par-yield moves of September 2024 to September 2026, unhedged and
with the minimum-variance two-, twenty- and thirty-year hedge (estimated on
2016--2026). Source: US Treasury par yields; the chapter's
tutorial.}\label{fig:rc:rates-risk:pnl}
\end{omfigure}

\begin{remark}[What the hedge assumes]\label{rem:rc:rates-risk:assumes}
The hedge is only as good as its covariance: ten years of Treasury moves
include the zero-rate years and the 2022 hikes, and swap rates move with
Treasury yields plus swap spreads, which this model ignores. A hedge estimated
on one regime and run in another leaves more than 4\%; stress tests
(\cref{ch:rc:stress-testing-and-scenarios}) exist for that.
\end{remark}

\section{Gamma and cross-gamma}

The ladder is a first derivative; for large moves the second derivatives
matter.

\begin{definition}[Cross-gamma]\label{def:rc:rates-risk:crossgamma}
The \emph{cross-gamma}\index{cross-gamma} of a position between inputs $k$ and
$l$ is $\Gamma_{kl} = \partial^2V/\partial q_k\partial q_l$ for $k\ne l$, per
square basis point; with the diagonal terms it forms the gamma matrix, and the
second-order P\&L is $\sum_k\Delta_k\delta q_k +
\tfrac12\sum_{k,l}\Gamma_{kl}\delta q_k\delta q_l$.
\end{definition}

\begin{example}[The second-order term of the hook]\label{ex:rc:rates-risk:gamma}
On the move of 21 October 2022 the book's ladder gives a first-order loss of
exactly USD~3\,000\,000 (by construction); full revaluation gives
$-3\,077\,970$. The gamma matrix, computed by central differences, gives
$\tfrac12\,\delta q^\top\Gamma\,\delta q = -78\,189$, which accounts for the
difference to within USD~220. The book's parallel gamma is $-377$ dollars per
square basis point: it is short convexity, mostly through its thirty-year
payers.
\end{example}

For a swap book \omterm{def:rc:rates-risk:crossgamma}{cross-gamma} is small and mostly comes from discounting; for
options, and for products whose payoff depends on two rates (chapter~6), it is
first order.

\section{Tutorial: ladder, factors and hedge}\label{tut:rc:rates-risk:tutorial}

\begin{tutorial}
\textbf{Goal.} Measure the book's risk in quote, zero and key-rate
coordinates, find how the Treasury curve moves, and hedge the book with three
swaps. \textbf{End state:}
\cref{fig:rc:rates-risk:ladder,fig:rc:rates-risk:pca,fig:rc:rates-risk:pnl} and the
numbers of \cref{ex:rc:rates-risk:hedge}.
\begin{tutsteps}
\item \textbf{The book}: \texttt{balanced\_book()} returns 202 swaps with zero
parallel DV01.
\item \textbf{Ladders}: \texttt{zero\_and\_key\_ladders(book)}; check that
\texttt{par\_to\_zero} of the par ladder equals the zero ladder.
\omcode{code/firm/ratesrisk/firm_ratesrisk.py}{19}{27}{The par ladder: bump each quote, recalibrate, reprice.}\label{lst:rc:rates-risk:ladder}
\item \textbf{Factors}: \texttt{components()} on the Treasury file; the shares
should read 85.1, 11.0 and 2.1\%.
\item \textbf{Hedge}: \texttt{hedge(book)} and \texttt{hedge(book, (1, 5, 7))}
for the two-, ten- and thirty-year trio; \texttt{fig\_rc\_ratesrisk.py} writes
the charts.
\end{tutsteps}
\textbf{What to change next.} Estimate the covariance on 2016--2020 only and
measure the hedged variance on 2021--2026; replace flat forwards by monotone
convex and compare the ladders.
\end{tutorial}

\section{Build: the risk-transformation toolkit}\label{bld:rc:rates-risk:ratesrisk}

\begin{build}
\bfield{Purpose} The rates-risk layer of the miniature firm: every rates
position's ladder, in any coordinates, and the hedge that neutralises it; the
risk engine of \cref{ch:rc:build-a-risk-engine} calls it.
\bfield{Interface} \texttt{par\_ladder(spot, instruments, kind, pv)},
\texttt{zero\_ladder(curve, pv)}, \texttt{par\_to\_zero}, \texttt{zero\_to\_par},
\texttt{TentBumped}, \texttt{key\_rate\_ladder}, \texttt{pca(changes)},
\texttt{min\_variance\_hedge(g, H, cov)}, \texttt{cross\_gamma}.
\bfield{Rules} Sensitivities per basis point from small bumps scaled up
(chapter~1's nonlinearity); PCA signs fixed (level positive, slope rising with
maturity); no global state.
\bfield{Acceptance tests} \texttt{code/firm/ratesrisk/tests/}: the Jacobian maps
par ladders to zero ladders under two interpolations; key-rate sensitivities
sum to the parallel zero sensitivity; PCA recovers a planted two-factor
structure; hedging a book with itself leaves nothing; the gamma matrix is
symmetric and explains the second-order P\&L of a thirty-year swap.
\bfield{Stretch} Forward-rate buckets; an analytic Jacobian from
chapter~1's solver; a hedge with transaction costs (a penalty on $|h|$).
\end{build}

\begin{omsources}
\item R.\ Litterman and J.\ Scheinkman, ``Common factors affecting bond
returns'', \emph{Journal of Fixed Income} 1(1), 1991.
\item T.\ S.\ Y.\ Ho, ``Key rate durations: measures of interest rate risks'',
\emph{Journal of Fixed Income} 2(2), 1992.
\item US Department of the Treasury, \emph{Daily Treasury Par Yield Curve Rates}.
\item L.\ Andersen and V.\ Piterbarg, \emph{Interest Rate Modeling}, volume 3,
Atlantic Financial Press, 2010, chapter 26 (vega and delta hedging of rates books).
\end{omsources}

\section{Exercises}

\begin{exercise}[$\star$]\label{exo:rc:rates-risk:1}
A book's ladder is $+50$, $-120$, $+80$ thousand dollars per basis point on
the two-, ten- and thirty-year quotes. Give its parallel DV01 and its P\&L if
the three quotes move by $-5$, $+2$ and $+6$ basis points.
\end{exercise}

\begin{exercise}[$\star$]\label{exo:rc:rates-risk:2}
Why does the book's par ladder sum to zero while its zero ladder sums to
USD~11\,408?
\end{exercise}

\begin{exercise}[$\star$]\label{exo:rc:rates-risk:3}
Using \cref{ex:rc:rates-risk:pca}, what share of daily variance is left after
the first three components, and what is a one-standard-deviation daily move of
the \omterm{def:rc:rates-risk:factors}{level factor}, in basis points of score?
\end{exercise}

\begin{exercise}[$\star\star$]\label{exo:rc:rates-risk:4}
Compute the ten-year payer's zero-ladder sum from \cref{ex:rc:rates-risk:twoladders}
and compare it with its par ladder. Which is larger, and why?
\end{exercise}

\begin{exercise}[$\star\star$]\label{exo:rc:rates-risk:5}
Why is a key-rate ladder independent of the curve's interpolation while a par
ladder is not? Which one would you send to a regulator, and which one to a
trader who will hedge with par swaps?
\end{exercise}

\begin{exercise}[$\star\star$]\label{exo:rc:rates-risk:6}
Read \cref{fig:rc:rates-risk:ladder}. After the minimum-variance hedge the
seven- and ten-year buckets are unchanged at $+250$ and $-280$ thousand. Why
did the hedge leave them, and what do they have in common?
\end{exercise}

\begin{exercise}[$\star\star\star$]\label{exo:rc:rates-risk:7}
\emph{Coding.} Find the best single-swap hedge and the best four-swap hedge of
the book, and the variance each leaves.
\end{exercise}

\begin{exercise}[$\star\star\star$]\label{exo:rc:rates-risk:8}
\emph{Find the flaw.} ``We neutralise the level, slope and \omterm{def:rc:rates-risk:factors}{curvature factors}
every evening, so we carry no curve risk.'' Correct it with the book's numbers.
\end{exercise}

\section{Problem: The Twisted Day}

\begin{problem}[{Weekend problem --- a DV01-flat book on 21 October 2022}]\label{pb:rc:rates-risk:1}
The chapter's book (202 swaps, parallel DV01 zero) is replayed on the par-yield
moves of 21 October 2022: $-8$, $-13$, $-14$, $-11$, $-8$, $-3$, $+7$, $+9$
basis points from one to thirty years.

\medskip\noindent\textbf{Part I --- The loss.}
\begin{enumerate}
\item Give the first-order P\&L from the ladder.
\item Give the P\&L by full revaluation, and the difference.
\item Which two buckets contribute most, with their signs?
\item Was the move a steepening or a flattening, and which way was the book
positioned?
\item Why did the parallel DV01 give no warning?
\end{enumerate}

\medskip\noindent\textbf{Part II --- Factors.}
\begin{enumerate}[resume]
\item Give the move's score on the \omterm{def:rc:rates-risk:factors}{slope factor}.
\item Give the \omterm{def:rc:rates-risk:factors}{slope factor}'s contribution to the loss.
\item How many standard deviations was that slope move?
\item The higher components contribute $+0.72$, $+0.69$ and $-1.48$ million. What do
they represent?
\item Would a factor model with three factors have predicted the loss?
\end{enumerate}

\medskip\noindent\textbf{Part III --- Hedging.}
\begin{enumerate}[resume]
\item Give the unhedged daily standard deviation of P\&L.
\item Give the residual share with two-, ten- and thirty-year swaps.
\item Give the best three swaps, their notionals and directions.
\item Give the residual daily standard deviation with them.
\item Give the first-order P\&L of the hedged book on 21 October 2022.
\end{enumerate}

\medskip\noindent\textbf{Part IV --- Judgement.}
\begin{enumerate}[resume]
\item Why not simply hedge every bucket to zero?
\item How often should the covariance be re-estimated, and on what window?
\item What would you add to the daily risk report?
\item State the \emph{named result}: the loss on the twisted day, and the three
swaps that remove most of the risk with the variance they leave.
\item In one sentence: why is DV01 not a risk measure?
\end{enumerate}
\end{problem}

\section{Interview questions}

\begin{interviewq}[$\star$ \iqroles{trader, risk}]\label{iq:rc:rates-risk:1}
Your book has zero DV01. Are you hedged?
\end{interviewq}

\begin{interviewq}[$\star\star$ \iqroles{researcher, developer}]\label{iq:rc:rates-risk:2}
How do you convert risk to par-swap quotes into risk to zero rates, and back?
\end{interviewq}

\begin{interviewq}[$\star\star$ \iqroles{researcher, trader}]\label{iq:rc:rates-risk:3}
What are the first three principal components of yield-curve changes, and how
much do they explain?
\end{interviewq}

\begin{interviewq}[$\star\star$ \iqroles{trader}]\label{iq:rc:rates-risk:4}
You must hedge a book with only three swaps. How do you choose them?
\end{interviewq}

\begin{interviewq}[$\star\star\star$ \iqroles{risk, researcher}]\label{iq:rc:rates-risk:5}
Your first-order P\&L explain misses by 3\% on a big-move day. What is missing
and how do you compute it?
\end{interviewq}

\begin{interviewq}[$\star\star\star$ \iqroles{researcher}]\label{iq:rc:rates-risk:6}
When would you not trust a PCA-based hedge?
\end{interviewq}
